% Blueprint for the formalization of
%   H. Hironaka, Resolution of singularities of an algebraic variety over a field of
%   characteristic zero I, II.  Ann. of Math. 79 (1964), 109-203 and 205-326.
%
% All nodes are generated from the Lean sources by LeanArchitect (`lake build :blueprint`).
% This file only fixes the chapter structure, which follows the paper.

\input{../../.lake/build/blueprint/library/Hironaka}

\chapter*{Scope and conventions}

This blueprint certifies Hironaka's 1964 proof \emph{as written}.  Every node is labelled
by the paper's own numbering (\texttt{def:I.3-I} = Definition 3 (I) of Chapter I, etc.).
Material that Hironaka defers to forthcoming papers (Ch.~0 \S1, \S6,
\S7 on analytic spaces and schematization; Ch.~I \S4 ``details left to the reader'') is
out of scope.

Conventions of the formalization (see the module headers for details): subschemes are
coherent sheaves of ideals; containment is reverse inclusion of ideal sheaves; dimensions
are Krull dimensions in $\mathbb N\cup\{\infty\}\cup\{-\infty\}$; the fundamental theorems
are indexed by $n\in\mathbb N$, the case $n=-1$ being stated separately.

\chapter{Compatibility layer}

\section{Regular local rings and the order function}
\inputleannode{def:rsop}
\inputleannode{lem:regular-local-ring-iff-rsop}
\inputleannode{def:order-nu}
\inputleannode{lem:order-nu-basic}

\section{Regular local rings are domains; primes with regular quotient}

Facts that the paper uses without comment and that Mathlib does not provide: a regular local ring
is a domain, and a prime $P$ with $R/P$ regular is generated by part of a regular system of
parameters, with $P^\nu$ $P$-primary for all $\nu$.  All are proved.

\inputleanmodule{Hironaka.Compat.RegularQuotient}

\section{The associated graded ring of a regular local ring}

``We can identify $\operatorname{gr}_P(R)$ with the polynomial ring $\bar R[Z_1,\dots,Z_r]$''
(Ch.~II, p.~191; also pp.~183, 225): for a prime $P$ with $R/P$ regular this is proved here, via
regular sequences being quasi-regular (Matsumura, Theorems 14.4 and 16.2).

\inputleanmodule{Hironaka.Compat.RegularGraded}

Footnote~39 of Ch.~II (p.~183): for a prime $N$ of a regular local ring $A$ with $A/N$ regular, a system of elements of $N$
generates $N$ if and only if its images in $\mathfrak M/\mathfrak M^2$ span $(N+\mathfrak M^2)/\mathfrak M^2$; and the regular
system of parameters adapted to $P\subseteq Q$ of footnote~41 (p.~191), built from it as printed.

\inputleanmodule{Hironaka.Compat.RegularQuotientGenerators}

\section{Localizations of regular local rings (Serre)}
\inputleanmodule{Hironaka.Compat.RegularLocalization}

\section{Completions of regular local rings}
\inputleanmodule{Hironaka.Compat.AdicCompletionRegular}

\section{Helpers for the class $\mathfrak B$ (Ch.~I, p.~161)}
The regular locus of $\operatorname{Spec}A$, its transport along ring isomorphisms and along the
canonical map to the completion of a complete local ring; used for (a) and (b) of p.~161
(\cref{prop:class-A-a,prop:class-A-b}).

\inputleanmodule{Hironaka.Compat.HironakaBaseStability}

\section{Property (c) of the class $\mathfrak B$ (Ch.~I, p.~161, footnotes~30, 31)}
Footnote~31: ``Reduce the proof to the case in which S is complete, and then apply the jacobian
criterion.'' The cited inputs are Cohen's structure theorem and the jacobian criterion (footnote~30;
Nagata [14]), and two facts of local algebra cited
through footnote~3 of the Introduction (descent of regularity along flat local homomorphisms, and
flatness of the induced map of completions). Both steps of footnote~31 (the reduction to complete $S$ and the complete case) are proved, modulo these cited leaves and Serre's theorem (footnote~9).

\inputleanmodule{Hironaka.Compat.FlatLocalRegularity}
\inputleanmodule{Hironaka.Compat.PowerSeriesPolynomialRing}
\inputleanmodule{Hironaka.Compat.JacobianCriterion}
\inputleanmodule{Hironaka.Compat.PowerSeriesPresentation}
\inputleanmodule{Hironaka.Compat.QuotientCompletion}
\inputleanmodule{Hironaka.Compat.RegularQuotientIndependent}
\inputleanmodule{Hironaka.Compat.IndepModSqHeight}
\inputleanmodule{Hironaka.Compat.JacobianMinors}
\inputleanmodule{Hironaka.Compat.DerivationExtension}
\inputleanmodule{Hironaka.Compat.MvPowerSeriesRegular}
\inputleanmodule{Hironaka.Compat.PowerSeriesPolynomialRingRegular}
\inputleanmodule{Hironaka.Compat.PowerSeriesPolynomialRingExtension}
\inputleanmodule{Hironaka.Compat.ClassBCompleteCore}
\inputleanmodule{Hironaka.Compat.RegularLocusFlatDescent}
\inputleanmodule{Hironaka.Compat.ClassBCompleteTransport}
\inputleanmodule{Hironaka.Compat.ClassBCompleteKernel}
\inputleanmodule{Hironaka.Compat.ClassBCompleteCase}
\inputleanmodule{Hironaka.Compat.RingPushoutSquares}
\inputleanmodule{Hironaka.Compat.RegularLocusPushout}
\inputleanmodule{Hironaka.Compat.ClassBReductionFibre}
\inputleanmodule{Hironaka.Compat.ClassBReductionDescent}
\inputleanmodule{Hironaka.Compat.ClassBLocalization}

\section{The local criterion of flatness (Ch.~II \S1, Lemma~2)}

Lemma~2 of Ch.~II \S1 (p.~182) is cited from Grothendieck; Mathlib has neither the local
criterion of flatness nor a long exact sequence for $\operatorname{Tor}$.  The condition
$\operatorname{Tor}_1^A(M,A/I)=0$ is identified with injectivity of $M\otimes_AI\to M$, and the
implications (ii)$\Rightarrow$(i) and (iii)$\Rightarrow$(ii) are proved via Artin--Rees and
Krull's intersection theorem.

\inputleanmodule{Hironaka.Compat.TorOne}
\inputleanmodule{Hironaka.Compat.LocalFlatness}

For the printed proof of Lemma~4 of Ch.~II \S1 (p.~188): the base change $A'\otimes_AB$, $A'\otimes_AM$ of a graded algebra
and module; $\operatorname{Tor}_1^A(M,A/P)$, in the kernel model, as a graded $B$-module of finite type; and its
compatibility with localization.

\inputleanmodule{Hironaka.Compat.GradedBaseChange}
\inputleanmodule{Hironaka.Compat.TorOneGraded}
\inputleanmodule{Hironaka.Compat.TorOneLocalization}

\section{Flatness of RSOP-preserving homomorphisms (Ch.~III \S2, p.~207)}

The proof of the Corollary to Lemma~2 of Ch.~III \S2 asserts without proof that a local
homomorphism of regular local rings mapping a regular system of parameters onto one is flat
(p.~207; used again on p.~216).  This gap is filled here: induction on the number of
parameters, passing to the quotients by one parameter, and the local criterion of flatness
above (Matsumura, Theorem~23.1, specialised).

\inputleanmodule{Hironaka.Compat.RSOPFlat}

\section{Generic freeness (Ch.~II \S1, Lemma~3)}

Lemma~3 of Ch.~II \S1 (p.~188) is cited from Grothendieck; Mathlib does not have it.  It is
proved by induction on the number of generators of the algebra, filtering the module by degree
in the last generator.

\inputleanmodule{Hironaka.Compat.GenericFreenessAux}
\inputleanmodule{Hironaka.Compat.GenericFreeness}

\section{Associated graded algebras (Ch.~II \S1--\S2, pp.~180, 189)}

The graded algebra $\operatorname{gr}_I(A)=\bigoplus_m I^m/I^{m+1}$ of p.~180 is realized as the
quotient of the Rees algebra $\bigoplus_m I^mX^m$ by the ideal generated by $I$; the order
$\nu_I(f)$, initial forms and the graded ideal $\operatorname{gr}_I(J,A)$ are those of p.~189, and
Lemma~5 of Ch.~II \S2 lives here because it is part of that setup.

\inputleanmodule{Hironaka.Compat.Graded}

\section{Stalks of subschemes and the simple locus}
\inputleanmodule{Hironaka.Compat.SubschemeStalk}

\section{Stalks of sheaves of ideals}
\inputleanmodule{Hironaka.Compat.IdealSheafStalk}

\section{Parts of regular systems of parameters}
\inputleanmodule{Hironaka.Compat.RSOPProducts}

\section{Affine charts of a monoidal transformation}

The paper describes the local rings of a monoidal transformation by charts
(``$\mathcal O_{x'}$ is a localization of $\mathcal O_x[z_{r+1}/z_t,\dots,z_s/z_t]$'', p.~169;
``we see that $R'$ is a monoidal transform of $R$'', p.~214).  Since monoidal transformations
are given by their universal property (Ch.~0 \S2), the charts are constructed here from that
property.

\inputleanmodule{Hironaka.Compat.BlowupCharts}

The monoidal transformation is an isomorphism over $X-B$ (p.~167), with no further hypothesis.
\inputleanmodule{Hironaka.Compat.BlowupLocalIso}

\section{Flat base change of monoidal transformations}

Ch.~IV \S1 uses that localization commutes with monoidal transformations (``$h_S$ is the
monoidal transformation of $X''_S$ with center $B''_S$''); since monoidal transformations are
given by their universal property, this is proved here from that property.

\inputleanmodule{Hironaka.Compat.BlowupBaseChange}

\section{Irreducible components through a point and stalks}
\inputleanmodule{Hironaka.Compat.ComponentStalk}

\chapter{Chapter 0: preliminaries}

\section{The class of base rings and algebraic schemes (Ch.~0 \S1, Ch.~I)}
\inputleannode{def:class-A}
\inputleannode{prop:class-A-iii}
\inputleannode{lem:class-A-complete}
\inputleannode{prop:class-A-a}
\inputleannode{prop:class-A-b}
\inputleannode{prop:class-A-c}
\inputleannode{prop:class-A-abc}
\inputleannode{def:algebraic-scheme}
\inputleannode{def:I.0-dimension-at}
\inputleannode{def:I.0-dimension}
\inputleannode{lem:I.0-dimension-topological}
\inputleannode{def:I.0-codimension-at}
\inputleannode{def:I.0-codimension}
\inputleannode{def:I.0-everywhere-codimension}

The codimension of Ch.~I, p.~163, as printed: the least $\dim_x(V)$ over the points $x$ of $W$.
\inputleanmodule{Hironaka.Compat.CodimDefinition}

\section{Invertible ideals and codimension one (Ch.~0 \S5, Definition~2)}
Definition~2 (p.~141) asserts, without proof or citation, that for a reduced subscheme $E$ of a
non-singular algebraic scheme the sheaf of ideals of $E$ is invertible ``or, equivalently,'' $E$ is
everywhere of codimension one. The equivalence is proved here (gap filled), from the factoriality
of regular local rings, which belongs to the theory of local rings cited in footnote~3 of the
Introduction (p.~116).

\inputleanmodule{Hironaka.Compat.RegularLocalRingUFD}
\inputleanmodule{Hironaka.Compat.CodimOneInvertible}

\section{Generalities on blowing-up (Ch.~0 \S2, Ch.~I \S1)}
\inputleannode{def:total-transform}
\inputleannode{def:monoidal-transformation}
\inputleannode{thm:monoidal-transformation-exists}
\inputleannode{lem:monoidal-transformation-unique}
\inputleannode{lem:monoidal-transformation-proper}
\inputleannode{lem:monoidal-transformation-empty-center}
\inputleannode{def:strict-transform}
\inputleannode{lem:strict-transform-saturation}
\inputleannode{lem:strict-transform-is-blowup}
\inputleannode{lem:strict-transform-is-blowup-unique}
\inputleannode{def:succession}
\inputleannode{lem:succession-proper}

\section{Local structure of a monoidal transformation (Ch.~0 \S5)}

The non-singularity of $X'$ and of $f^{-1}(B)$ is proved locally from Ch.~III \S2 (a monoidal
transform of a regular local ring with regular center is regular); the Lean module is placed
after Ch.~III \S2 for this reason.

\inputleanmodule{Hironaka.Ch0.BlowupLocal}

\section{Hilbert--Samuel functions and multiplicity}

Hironaka uses the multiplicity $e(\mathcal O_{X,x})$ of the local rings of $X$ in the sense of
Samuel (Zariski--Samuel, Vol.~II, Ch.~VIII), without redefining it.  We record the definition
and the one fact about it that the paper relies on, which is cited.

\inputleannode{def:hilbert-samuel-function}
\inputleannode{thm:hilbert-samuel}
\inputleannode{def:multiplicity}
\inputleannode{lem:0.4-multiplicity-hypersurface}
\inputleannode{def:0.4-multiplicity-at}
\inputleannode{def:equimultiple-along}

\section{Normal flatness (Ch.~0 \S4, Definition 1)}
\inputleannode{def:normally-flat-ideal}
\inputleannode{lem:normally-flat-iff-flat}
\inputleannode{lem:0.1-normally-flat-iff-gr-flat}
\inputleannode{def:0.1-normally-flat-at}
\inputleannode{def:0.1-normally-flat}

\section{Normal flatness and equimultiplicity (Ch.~0 \S4 (I), (II))}

Hironaka proves (I) and (II) of Ch.~0 \S4 in Ch.~II (footnotes 18, 19); the Lean module is
placed after Ch.~II for this reason.

\inputleanmodule{Hironaka.Ch0.NormalFlatnessCriteria}

The converse of (I) of Ch.~0 \S4 (p.~136) with its printed hypotheses ($W$ possibly reducible, the ideal of $V$ locally
principal, zero allowed).
\inputleanmodule{Hironaka.Compat.LocallyPrincipalConverse}

\section{Normal crossings (Ch.~0 \S5, Definition 2)}
\inputleannode{def:0.2-normal-crossings-with-at}
\inputleannode{def:0.2-normal-crossings-at}
\inputleannode{def:0.2-normal-crossings}
\inputleannode{def:0.2-normal-crossings-with}

\section{Openness of the normal crossing locus (Ch.~0 \S5; Ch.~IV \S1, Lemma~1)}
\inputleanmodule{Hironaka.Ch0.NormalCrossingsOpen}

\chapter{Chapter I: an inductive formulation of resolution problems}

\section{Resolution data and permissible transformations (\S1)}
\inputleannode{def:I.3-I}
\inputleannode{def:I.4-I}
\inputleannode{def:I.3-II}
\inputleannode{def:I.4-II}
\inputleannode{def:I.S}
\inputleannode{def:I.nu}
\inputleannode{def:I.S-nu}
\inputleannode{def:I.S-star}
\inputleannode{def:I.5-I.1}
\inputleannode{def:I.5-I.2}
\inputleannode{def:I.5-II.1}
\inputleannode{def:I.6-I}
\inputleannode{def:I.6-I-closed}
\inputleannode{def:I.6-I-open}
\inputleannode{def:I.6-II}
\inputleannode{def:I.6-II-closed}
\inputleannode{compat:blowup-stalk-chart-printed}
\inputleannode{lem:I.nu-attained}
\inputleannode{rem:I.1-dim}
\inputleannode{rem:I.1-dim-empty}
\inputleannode{rem:I.1-dim-typeII}
\inputleannode{rem:I.1-dim-typeI}
\inputleannode{def:I.7-I}
\inputleannode{def:I.7-II}
\inputleannode{def:I.7-II-as-printed}
\inputleannode{rem:I.1-as-printed-false}
\inputleannode{lem:I.7-II-exponent-nonneg}
\inputleannode{rem:I.1-transform-is-datum-I}
\inputleannode{rem:I.1-dim-W}
\inputleannode{rem:I.1-transform-is-datum-II}
\inputleannode{lem:I.7-II-divisible}
\inputleannode{lem:0.5-weak-transform-order}
\inputleannode{lem:I.7-II-weak-transform}
\inputleannode{rem:I.1-transform-is-datum-II-quotient}
\inputleannode{def:I.8}
\inputleannode{lem:I.8-transform-is-datum}

\section{The implication diagram (\S2)}

This section is the purely combinatorial core of Ch.~I \S2, formalized without Mathlib: an
abstract family of propositions indexed like the fundamental theorems, the four implications
(A)--(D) of pp.~171--172 as properties of such a family, and the theorem that (A)--(D)
together with the trivial base cases yield every $\mathrm I_1^{N,n}$, $\mathrm I_2^{N,n}$,
$\mathrm{II}_1^N$, $\mathrm{II}_2^N$.  This is the one step of \S2 that Hironaka leaves to the
reader (``this is easily done'').  The instances of (A)--(D) for the actual fundamental theorems, which are
the content of Chapters~II--IV, are stated in the next section.

\inputleanmodule{Hironaka.Ch1.InductionScheme}

\section{The fundamental theorems (\S2)}
\inputleannode{thm:I.2-I1}
\inputleannode{thm:I.2-I2}
\inputleannode{thm:I.2-II1}
\inputleannode{thm:I.2-II2}
\inputleannode{def:I.2-fundamental-family}
\inputleannode{lem:I.2-trivial-I-n0}
\inputleannode{lem:I.2-trivial-II-N0}
\inputleannode{thm:I.2-I1-minus-one}
\inputleannode{thm:I.2-I2-minus-one}
\inputleannode{lem:I.2-trivial-I-minus-one}

\chapter{Chapter II: normal flatness along subschemes}

The statements of this chapter are transcribed from pp.~180--203.  Flatness of $\operatorname{gr}_P(O)$
over $O/P$ is expressed through Definition~1 of Ch.~0 (freeness of the homogeneous parts,
cf.\ p.~136).

\section{Flatness of the associated graded algebra (\S1)}
The graded module $\gr_I(M)$ of p.~180 as a graded $\gr_I(A)$-module of finite type, and the canonical map $\varphi$ as
one map.
\inputleanmodule{Hironaka.Compat.GradedModule}

\inputleanmodule{Hironaka.Ch2.Flatness}

Proposition~1 (1) (p.~184) as printed: $\psi=\alpha\otimes\beta^*$ as one isomorphism of graded $K$-algebras
$\gr^0_{\bar m}(\gr_P(O))\otimes_K\gr_{\bar m}(\bar O)\to\gr_m(O)$.
\inputleanmodule{Hironaka.Ch2.PsiAlgebra}

\section{The associated graded ideal (\S2)}
\inputleanmodule{Hironaka.Ch2.GradedIdealBase}

Lemma~8 (pp.~193--194) by the printed diagram: $\psi'$ is an isomorphism if and only if $\psi''$ is surjective, and the
image of $\psi''$ is generated by $\gamma(\gr_P(J,R))$.
\inputleanmodule{Hironaka.Ch2.Lemma8Diagram}

The exact sequences of graded modules in the proof of Theorem~2 (pp.~195--196): ``it suffices that
$\operatorname{Ker}(\varepsilon')=\operatorname{Im}(\delta')$'' and the chase giving the injectivity of $\varphi$.
\inputleanmodule{Hironaka.Ch2.GradedModuleExact}

\inputleanmodule{Hironaka.Ch2.GradedIdeal}

Corollary~2 of Theorem~2 (p.~197) with the printed hypothesis ``locally everywhere generated by a single non-zero element''
(on a non-singular scheme this is equivalent to invertibility; gap filled).
\inputleanmodule{Hironaka.Compat.InvertibleNonzeroGenerator}

\inputleanmodule{Hironaka.Ch2.GradedIdealCorollary2}

\section{Normal flatness along an ascending chain of subschemes (\S3)}
\inputleanmodule{Hironaka.Ch2.AscendingChain}

\section{Corollary 1 of Proposition 1 at a point (used in Ch.~IV \S2)}
\inputleanmodule{Hironaka.Ch2.PointwiseFlatness}

\chapter{Chapter III: effects of permissible monoidal transformations on singularities}

The definitions shared by the sections of this chapter (the characters $\nu^{(i)}(J)$ and
$\nu^*(J)$, standard bases, monoidal transforms of a regular local ring, residually
algebraic/rational/separable local homomorphisms, and the strict and weak transforms of an
ideal, pp.~205--214) are collected first; each section then follows the paper.

\section{Shared definitions (\S1--\S2, pp.~205--214)}
\inputleanmodule{Hironaka.Ch3.Basic}

\section{The characters $\nu^*(J)$ and standard bases (\S1)}
\inputleanmodule{Hironaka.Ch3.NuStar}

\section{Tensor products of field extensions (used in \S2, Lemma~4)}
\inputleanmodule{Hironaka.Compat.FieldTensor}

\section{The strict and weak transforms of a local ideal (\S2)}
Associated primes of a quotient along a surjection (for the correspondences in the proof of Lemma~4, p.~211).
\inputleanmodule{Hironaka.Compat.AssociatedPrimes}

\inputleannode{lem:III.2-regular-local-domain}
\inputleannode{lem:III.2-chart}
\inputleannode{lem:III.2-monoidal-transform-index}
\inputleannode{lem:III.2-monoidal-transform-regular}
\inputleannode{lem:III.2-monoidal-transform-principal}
\inputleannode{lem:III.2-dim-rsop}
\inputleannode{lem:III.2-dim-maximal}
\inputleannode{lem:III.2-dim}
\inputleannode{lem:III.2-dim-finite}
\inputleannode{lem:III.2-residually-rational-rsop}
\inputleannode{def:III.2-canonical}
\inputleannode{lem:III.4-minimal-base}
\inputleannode{lem:III.4-associated-primes}
\inputleannode{lem:III.4-q-correspondence}
\inputleannode{lem:III.4-count}
\inputleannode{lem:III.4}
\inputleannode{rem:III.2-1}
\inputleannode{def:III.2-weak-transform-sheaf}
\inputleannode{rem:III.2-2}
\inputleannode{def:III.2-strict-transform-sheaf}
\inputleannode{rem:III.2-3}
\inputleannode{rem:III.2-4}
\inputleannode{lem:III.5}
\inputleannode{cor:III.lem5}
\inputleannode{lem:III.6}

The associated primes of $\bar{\mathfrak m}'\tilde k[X]$ are isolated, for any extension $\tilde k$ of $k$ (used in the proof of
Lemma~4, p.~211).
\inputleanmodule{Hironaka.Compat.EmbeddedPrimes}

Remark~0 (p.~213): the transform does not determine the center.
\inputleanmodule{Hironaka.Ch3.TransformsRemark0}

Existence of the $\tilde R'$ of Lemma~4 (1) from the associated primes of $M'\tilde A'$, as printed (p.~211).
\inputleanmodule{Hironaka.Compat.CanonicalLiftExistence}

A monoidal transform written with any generator of $P$ as denominator, as a ring of fractions of
$R[x_1/x_i,\dots,x_r/x_i]$ (p.~209; used in Remark~1 of Ch.~I, p.~169).
\inputleanmodule{Hironaka.Compat.TransformRingOfFractions}

\section{The completion: Corollaries 1 and 2 of Lemma 4 (\S2)}
\inputleanmodule{Hironaka.Ch3.TransformsCompletion}

\section{Regular systems of parameters of a monoidal transform (used in Ch.~I, Remark~1)}
Formalization note: these statements belong to Ch.~I \S1, Remark~1 (p.~169), and are placed
after Ch.~III \S2, whose theory of monoidal transforms of local rings they use.
\inputleanmodule{Hironaka.Compat.TransformRSOP}

\section{Dimension of a monoidal transform (used in Ch.~I, Remark~1)}
\inputleanmodule{Hironaka.Compat.BlowupDim}

\section{Resolving a curve germ by point blow-ups (used in Ch.~III, Theorem~1)}
\inputleanmodule{Hironaka.Compat.CurveResolution}

\section{Irreducibility of the exceptional subscheme (used in Ch.~I, Remark~1)}
\inputleanmodule{Hironaka.Compat.BlowupExceptional}

\section{The dimension inequality for birational extensions}
\inputleanmodule{Hironaka.Compat.DimInequality}

\section{Dimension of strict transforms (used in Ch.~I, Remark~1)}
\inputleanmodule{Hironaka.Compat.StrictTransformDim}

\section{The upper semi-continuity of $\nu(J)$ (\S3)}
\inputleanmodule{Hironaka.Ch3.SemiContinuity}

\section{The characters $\tau^*(J)$ and regular systems of $\tau$-parameters (\S4)}
\inputleanmodule{Hironaka.Ch3.TauStar}

\section{The maximal local extension (proof of Theorem~3, \S5, p.~234)}
\inputleanmodule{Hironaka.Compat.MaximalLocalExtension}

\section{Effects of permissible monoidal transformations on $\nu^*$ and $\tau^*$ (\S5)}
\inputleanmodule{Hironaka.Ch3.PermissibleEffects.Lemma14}
\inputleanmodule{Hironaka.Ch3.PermissibleEffects.Lemmas15And16}
\inputleanmodule{Hironaka.Ch3.PermissibleEffects.Theorems3And4}

\section{A stability theorem of a standard base (\S6)}
\inputleanmodule{Hironaka.Ch3.Stability}

\section{Corollary 3 (c) of Theorem 5 (\S6)}
\inputleanmodule{Hironaka.Ch3.StabilityCor3c}

\section{Regular frames and normalized standard bases (\S7)}
\inputleanmodule{Hironaka.Ch3.Frames}

Theorem~7 (p.~253) along the printed route: reduction to the complete case and the partial frames
$S_a=S\{z_{\tau_a+1},\dots,z_r\}$ (the proof sits after the partial-frame infrastructure of \S10).
\inputleanmodule{Hironaka.Compat.FrameTheorem7}

\inputleanmodule{Hironaka.Ch3.FramePermissibleCenters}

\section{Effects of permissible monoidal transformations on a normalized standard base (\S8)}
\inputleanmodule{Hironaka.Compat.CoefficientField}

\inputleanmodule{Hironaka.Compat.TauParametersInPrime}

\inputleanmodule{Hironaka.Compat.FrameWithElement}

\inputleanmodule{Hironaka.Ch3.FrameTransforms}

\inputleanmodule{Hironaka.Compat.RegularClosedPoint}

\inputleanmodule{Hironaka.Compat.PermissibleCenterBlowdown}

\inputleanmodule{Hironaka.Compat.DimensionFacts}
\inputleanmodule{Hironaka.Compat.NormalFlatnessLocalization}

\section{The notion of $J$-stability (\S9); the existence of a $J$-stable regular $\tau$-frame and a $J$-stable standard base (\S10)}

Sequences of monoidal transforms are the $\mathbb N$-indexed sequences of \S5, a finite sequence
being continued by the trivial centers $P_\alpha=R_\alpha$.

\inputleanmodule{Hironaka.Ch3.JStability}
\inputleanmodule{Hironaka.Ch3.JStabilityLemma22}
\inputleanmodule{Hironaka.Ch3.JStableStandardBase}

\inputleanmodule{Hironaka.Ch3.TauFrameExistence}

\inputleanmodule{Hironaka.Compat.WeierstrassFrameStep}


\chapter{Chapter IV: the fundamental theorems and their proofs}

This chapter establishes the basic implications (A)--(D) of Ch.~I \S2 (pp.~274--326).  A
resolution datum of type $\mathfrak R^{N,n}_{\mathrm I}$ is written, as in the chapter's
introduction (pp.~274--275), in the form
$(\bigcup_{i=1}^{\alpha}E_i;V_1,\dots,V_\beta;W)$ with every $E_i$ a non-singular irreducible
subscheme of codimension one (the empty subscheme allowed, repetitions allowed); in the code
the boundary is the explicit family $E\colon\mathrm{Fin}\,\alpha\to\mathrm{Subscheme}\,X$.

\section{Localizing and lifting successions (used in Ch.~IV \S1)}
\inputleanmodule{Hironaka.Compat.LocalizeSuccession}

The remaining sentences of Remark~1 of Ch.~I (p.~169): the datum is resolved everywhere when $X=B$, and
$\dim W'=\dim W$ for closed or open restrictions.
\inputleanmodule{Hironaka.Compat.RemarkOneDim}

\section{Localizing and lifting successions for data of type II (used in Ch.~IV \S1)}
\inputleanmodule{Hironaka.Compat.LocalizeSuccessionII}

\section{Dimension of the localization of a succession (used in Ch.~IV \S1, Remark~2)}
\inputleanmodule{Hironaka.Compat.LocalizeDim}

\section{Localization of resolution data and resolution problems (\S1)}
\inputleanmodule{Hironaka.Ch4.Localization}

\section{Preparation on resolution data $(\mathfrak R^{N,n}_{\mathrm I},U)$ (\S2)}
\inputleanmodule{Hironaka.Ch4.Preparation.Lemmas3To5}
\inputleanmodule{Hironaka.Ch4.Preparation.Proposition5Lemmas}
\inputleanmodule{Hironaka.Ch4.Preparation.Proposition5}

\subsection*{Extending successions from an open subscheme (Proposition 5, step 7, p.~291)}

\inputleanmodule{Hironaka.Compat.ExtendSuccession}

\subsection*{The strict transform meets the exceptional divisor irreducibly (Proposition 5, step (3), p.~290)}

\inputleanmodule{Hironaka.Compat.StrictTransformExceptional}

\subsection*{Regularity of the strict transform along a regular center contained in it (\S3, p.~298)}

The ``we see'' of p.~298: $\bar f'_s$ is the monoidal transformation of $\bar X'_s$ with the non-singular center
$\bar B'_s$, so $\bar X'_{s+1}$ is non-singular and the exceptional divisor is non-singular of codimension one.

\inputleanmodule{Hironaka.Compat.StrictTransformRegular}

\subsection*{The order of a stable element along a permissible center (\S3, p.~302)}

``It follows that $g_{jk}(s)\in\hat Q_s^{\nu_{jk}}$'' (p.~302): Definition~11 of Ch.~III applied to a sequence whose first
center is $\hat Q_s$.

\inputleanmodule{Hironaka.Compat.StableOrder}

\subsection*{Termination of the process of Lemma~5: the order of $E$ along $G$ (\S2, pp.~288--291)}

The order of $E_1|_W$ along a bad component $G$, at the generic point of $G$, drops by one under the monoidal
transformation with center $G$; so the process of Lemma~5 used in step (3)* of Proposition~5 stops.

\inputleanmodule{Hironaka.Compat.BlowupOrderDrop}

\subsection*{Localized frames for Lemma (AB.5) (\S3, pp.~303--305)}

``If $T_0$ denotes the localization of $T$ with respect to $\hat P_0\cap T$, then $(T_0;z)$ is a regular frame of
$\hat S_0$'' (p.~303), and ``(b) of (2) follows from (3) and (4)'' (pp.~304--305).

\inputleanmodule{Hironaka.Compat.FrameLocalization}

\subsection*{The step of the frames of Lemma (AB.5) at a non-closed point (\S3, p.~305)}

``Finally (a) of (2) follows (b) of (2) \dots\ by Corollary 3 of Theorem 5, \S6, Ch.~\textsc{iii}, and by Theorem 5'' (p.~305):
Theorem~5 for the monoidal transform with center $\hat{\mathbf Q}_t$, then Corollary~3 (c) for its localization.

\inputleanmodule{Hironaka.Compat.FrameStepLocal}

\subsection*{Singularity of $W(s)$ at $x'_s$ in Lemma (AB.5) (ii) (\S3, p.~305)}

``If $W$ is singular at $x_0$, then the number of those indices $k$ with $\nu_{0k}=1$ is less than codimension of $W$ in $X$
\dots\ It is clear that the codimension of $W(s)$ \dots\ at $x'_s$ is at least equal to that of $W$ in $X$ at $x_0$'' (p.~305).

\inputleanmodule{Hironaka.Compat.SingularStandardBase}

\subsection*{Descent of permissible centers in Lemma (AB.5) (iii) (\S3, pp.~305--306)}

``It follows that $\hat x'_t$ is a simple point of $c_t^{-1}(W(t))$'' and ``Since $\hat{\mathbf S}_{t+1}$ (resp.\ $\hat{\mathbf S}_t$) is
faithfully flat over $\mathbf S_{t+1}$ (resp.\ $\mathbf S_t$), it follows by Lemma 9 of \S3, Ch.~\textsc{ii} that \dots'' (p.~305).

\inputleanmodule{Hironaka.Compat.PermissibleCenterDescent}

\subsection*{Descent of normal crossings along $\mathbf S_s\to\hat{\mathbf S}_s$ (\S3, Remark~4, p.~309; \S4, Remark~1, p.~322)}

``the canonical homomorphism $\mathbf S_s\to\hat{\mathbf S}_s$ transforms a regular system of parameters of $\mathbf S_s$ into such a
system of parameters of $\hat{\mathbf S}_s$ \dots\ Therefore, by the above property of the homomorphism $\mathbf S_s\to\hat{\mathbf
S}_s$, we can conclude that $\bigcup_{i=1}^{\alpha+s}E_i(s)$ has only normal crossings with $B_s$ at $x$'' (p.~309).

\inputleanmodule{Hironaka.Compat.NormalCrossingsDescent}

\section{Proofs of the implications (A) and (B) (\S3)}
\subsection*{Frames shared by \S3 and \S4 (Ch.~III \S\S8--9)}

\inputleanmodule{Hironaka.Compat.FrameCore}

\subsection*{Points of a monoidal transformation and its local rings (for Lemma (AB.2), p.~295)}
A monoidal transform of the local ring at a point is the local ring of a point of the monoidal transformation over it, and
a point over a given point is determined by its local ring (from the universal mapping property of p.~166); with the count
of Lemma~4~(1) of Ch.~III, the one-step bijection of fibres.

\inputleanmodule{Hironaka.Compat.BlowupStalkPoints}
\inputleanmodule{Hironaka.Compat.BlowupFibreBijection}

Monoidal transformations commute with flat base change (for ``one can show that $\hat X_s$ is canonically isomorphic to
the product $X_s\times_X\hat X_0$'', p.~295).
\inputleanmodule{Hironaka.Compat.BlowupFlatPullback}

\inputleanmodule{Hironaka.Ch4.ImplicationsAB.LemmasAB1ToAB3}
\inputleanmodule{Hironaka.Ch4.ImplicationsAB.LemmasAB4AndAB5}
\inputleanmodule{Hironaka.Ch4.ImplicationsAB.PreparedSuccession}

\subsection*{Correction of Lemma (AB.4): the componentwise construction (pp.~298--302)}

Lemma (AB.4) asserts that $\hat B_s=c_s^{-1}(B_s)$ is irreducible, which is false; the condition [g] (p.~298)
then requires containments that the paper proves only on the component of $\hat B_s$ through $\hat x$. These nodes
apply [g] componentwise on $\hat X_s$ (only the components of $\hat B_s$ on which the containments fail are
removed), and prove the corrected (AB.4) and its Corollary by the paper's own argument. The printed construction
above is kept as printed.

\inputleanmodule{Hironaka.Ch4.ComponentwiseConstruction}

\subsection*{The Weierstrass form of the frame (p.~297)}

The algebra behind ``by replacing $y_j$ by $y_j+c_jz$ with suitable $c_j\in k$ \dots\ we may
assume'' (generic coordinates and Weierstrass preparation).

\inputleanmodule{Hironaka.Compat.GenericWeierstrass}

\subsection*{Lemma (CD.3) as printed: non-reduced subschemes over $\hat x_0$ (p.~316)}

The Corollary of Lemma (AB.3) is stated for reduced $\hat B$; Lemma (CD.3) refers to it but speaks of
irreducible $\hat B$. These nodes prove the statement without reducedness.

\inputleanmodule{Hironaka.Ch4.CompletionFibre}

\subsection*{The frame of \S4: ring-level data (pp.~314--315)}

\inputleanmodule{Hironaka.Ch4.CDFrame}

\subsection*{Lemma (AB.6): completion of strict transforms and permissible centers (p.~306)}

\inputleanmodule{Hironaka.Ch4.CompletionTransforms}

\subsection*{Lemma (AB.6), Remarks 3--5 and Theorems A*, B*, A, B (\S3, pp.~293--310)}

These nodes depend on Lemma (AB.6), whose proof uses the completion lemmas above; they are placed
after them.

\inputleanmodule{Hironaka.Ch4.TheoremsAStarBStar}

\subsection*{The final step of Theorems A*, B*: extending the succession (\S3, p.~310)}

\inputleanmodule{Hironaka.Ch4.StageState}

\inputleanmodule{Hironaka.Ch4.AStarBStarAssembly}

\inputleanmodule{Hironaka.Ch4.AStarBStarLifting}

\subsection*{Theorems A*, B*, A, B and the implications (A), (B) (\S3, pp.~293, 310)}

\inputleanmodule{Hironaka.Ch4.TheoremsAB}

\subsection*{Flat base change of strict transforms; stalks along the frame (Ch.~IV \S\S3--4)}

\inputleanmodule{Hironaka.Compat.FlatStrictTransform}

\inputleanmodule{Hironaka.Compat.StalkFrameStep}

\section{Proofs of the implications (C) and (D) (\S4)}

The section ends the paper: combined with \S3 and the induction scheme of Ch.~I \S2, Theorems
A--D yield all the fundamental theorems (\cref{thm:I.2-fundamental}; the deduction is checked
in Lean, modulo the cited results it depends on, which are hypotheses).

\inputleanmodule{Hironaka.Ch4.ImplicationsCD.LemmasD1AndD2}
\inputleanmodule{Hironaka.Ch4.ImplicationsCD.LemmasCD3ToCD5}
\inputleanmodule{Hironaka.Ch4.ImplicationsCD.PreparedSuccession}

\subsection*{Repair of the second error, \S4 side: the emptied auxiliary datum and the shifted threshold (pp.~321--326)}

\inputleanmodule{Hironaka.Ch4.AuxiliaryDatumRepair}

\subsection*{Theorem C (pp.~311, 314--326)}

\inputleanmodule{Hironaka.Ch4.TheoremC}

\subsection*{The end of the paper (p.~326)}

\inputleanmodule{Hironaka.Ch4.Conclusion}

\subsection*{Main Theorem II$(N)$ of Ch.~I \S3 (pp.~176--177)}
Deduced, as printed, from the fundamental theorem $\mathrm{II}_2^N$ applied to the datum with $b=1$ and all $a_i=0$; the reduction
to irreducible $X$ (``as is easily seen'') interleaves the successions of the irreducible components (gap filled).

\inputleanmodule{Hironaka.Compat.ResolveComponents}

\inputleannode{thm:I.3-main-II}
\inputleannode{lem:I.3-main-II-properties}

\subsection*{Main Theorem I$(n)$ of Ch.~I \S3 (pp.~172--176)}
The ambient result of pp.~172--174 (fundamental theorem $\mathrm I_2^{N,n}$, the strengthening (1$'$) and the elimination
of pp.~173--174), then the induction on $n$ of pp.~174--176 through the completion at a generic point of $\bar S(r)$,
the embedding of footnote~37 and the smallest $C$ with $\hat C=c_r^{-1}(C)$. The steps the paper leaves unargued are
gaps filled.

\inputleannode{def:I.3-main-system}
\inputleannode{def:I.3-main-resolved}
\inputleannode{def:I.3-main-S}
\inputleannode{def:I.3-main-step-with}
\inputleannode{def:I.3-main-step}
\inputleannode{def:I.3-main-4-data}
\inputleannode{def:I.3-main-123-with}
\inputleannode{def:I.3-main-123}
\inputleannode{def:I.3-main-123-hat}
\inputleannode{def:I.3-main-4}
\inputleannode{def:I.3-main-I-conclusion}
\inputleannode{def:I.3-main-S-at}
\inputleannode{def:I.3-main-image}
\inputleannode{def:I.3-reach4}
\inputleannode{lem:I.3-reach4-of-iso}
\inputleannode{def:I.3-main-I-statement}
\inputleannode{def:I.3-main-I-ambient-statement}
\inputleannode{lem:I.3-S-closed}
\inputleannode{lem:I.3-resolved-generization}
\inputleannode{lem:I.3-S-nowhere-dense}
\inputleannode{lem:I.3-resolved-off-exceptional}
\inputleannode{lem:I.3-main-4-iff-S-empty}
\inputleannode{lem:I.3-S-transform-subset}
\inputleannode{lem:I.3-main-image-append}
\inputleannode{lem:I.3-main-image-closed}
\inputleannode{lem:I.3-main-I-improve}

\inputleanmodule{Hironaka.Compat.CompletionFibreGeneral}

\inputleanmodule{Hironaka.Ch4.MainIAmbient}

\inputleanmodule{Hironaka.Ch4.MainTheoremI}

\subsection*{The main theorems as announced in Ch.~0 (pp.~132--143)}
Main Theorem~I in the strong form (pp.~132--133), Main Theorem~I$^*$ (p.~138), Main Theorem~II (p.~142) and (B$^*$)
(p.~143), deduced from Main Theorems I$(n)$ and II$(N)$ of Ch.~I \S3. The paper prints no derivation (``We shall prove in
this work''); each statement is proved for algebraic schemes over a ring of the class $\mathcal B$ (p.~151) with the
printed form over a field of characteristic zero as a corollary.

\inputleanmodule{Hironaka.Compat.ResolutionBookkeeping}

A succession of monoidal transformations is one monoidal transformation (pp.~132--133): the existence of the ideal $J(m)$
with $f_0^{-1}(J(m))=f_0^{-1}(J_0)^mJ_1$, proved through the fractions of the chart rings.

\inputleanmodule{Hironaka.Compat.BlowupCompositionMain}

\inputleanmodule{Hironaka.Compat.BlowupSingleCenter}

Main Theorem~I as printed (p.~132): a single center $D$ whose points are the singular locus, and the weak form.

\inputleanmodule{Hironaka.Ch4.MainTheoremIZeroCenter}

The Corollary of Main Theorem~I (p.~132): a non-singular projective model of a finitely generated field over a field of
characteristic zero, with the lemma, unproved and uncited in the paper, that a monoidal transformation of a projective
scheme is projective (proved by comparing the charts of the blow-up with those of $\operatorname{Proj}$).

\inputleanmodule{Hironaka.Compat.ProjectiveIdealForms}

\inputleanmodule{Hironaka.Compat.BlowupProjective}

\inputleanmodule{Hironaka.Ch4.ProjectiveModel}

\inputleanmodule{Hironaka.Ch4.MainTheoremIZero}

\inputleanmodule{Hironaka.Ch4.MainTheoremIIZero}

Rational applications and Question~(E), the elimination of points of indeterminacy (pp.~138--143): the composition of a
rational application with a morphism with dense preimages, points of indeterminacy, proper rational applications (p.~139),
the fact (*) of p.~141, which the paper uses without proof (a cited leaf), the two facts on ideals not divisible by an
invertible ideal that the reduction uses tacitly (gaps filled), and the affirmative answer to (E) from Corollary~1.

\inputleannode{def:0.5-preimage-dense}
\inputleannode{def:0.5-rational-precomposition}

\inputleannode{def:0.5-point-of-indeterminacy}
\inputleannode{def:0.5-is-morphism}
\inputleannode{def:0.5-proper-application}

\inputleanmodule{Hironaka.Compat.BlowupPreimageDense}

\inputleannode{def:0.5-primitive-ideal}
\inputleannode{lem:0.5-primitive-decomposition}

\inputleanmodule{Hironaka.Compat.PrimitiveWeakTransform}

\inputleanmodule{Hironaka.Ch0.FactStar}

\inputleannode{def:0.5-eliminates-indeterminacy}
\inputleannode{thm:0.5-question-E}
\inputleannode{cor:0.5-question-E-integral}

Corollary~1 of Main Theorem~II (pp.~143--144): the trivialization of a system of coherent sheaves of ideals.

\inputleanmodule{Hironaka.Ch4.MainTheoremIICorollaries}

The consequences (1)--(3) of Corollary~1 (pp.~144--145): pairs $(X,\theta)$ of a model and an isomorphism of function fields,
domination, and the class $\mathfrak M(\mathbf L/\mathbf B)$; (1) dominating models obtained by monoidal transformations;
(2) the vanishing of $R^if_*\mathcal O_{X_1}$ for $i>0$ and $R^0f_*\mathcal O_{X_1}=\mathcal O_{X_2}$; (3) non-singular
projective models. The statements that these arguments use without proof are cited leaves: the single blow-up case of (2)
and the two sentences of (3) (``The existence of such $(X_2,\theta_2)$ is obvious'' and ``necessarily projective''). The
properness of the rational application between two proper pairs, which the arguments use tacitly, is proved, and so are
the Leray edge maps of a composite, which the argument of (2) uses tacitly (\cref{thm:0.5-leray-edge-maps}: the edge maps of
the Grothendieck spectral sequence for $f_*$ and $g_*$, as $f_*$ sends injective modules to flasque ones, which are
$g_*$-acyclic).

\inputleannode{def:0.5-nowhere-dense-centers}

\inputleanmodule{Hironaka.Compat.PairSuccession}

\inputleannode{def:0.5-pair}
\inputleannode{def:0.5-rational-application-of-pairs}
\inputleannode{def:0.5-dominates}
\inputleannode{def:0.5-class-U}
\inputleannode{def:0.5-obtained-by}

\inputleanmodule{Hironaka.Compat.FunctionFieldPairCited}

\inputleanmodule{Hironaka.Ch4.MainTheoremIIConsequence1}

\inputleanmodule{Hironaka.Compat.HigherDirectImageStructureSheaf}

\inputleanmodule{Hironaka.Ch4.MainTheoremIIConsequence2}

\inputleanmodule{Hironaka.Ch4.MainTheoremIIConsequence3}

Corollary~2 of Main Theorem~II (pp.~145--146): after monoidal transformations with nowhere dense centers, the pull-back of
a non-zero locally free coherent sheaf contains an invertible subsheaf with locally free quotient.

\inputleanmodule{Hironaka.Ch4.MainTheoremIICorollary2}

Corollary~3 of Main Theorem~II (p.~146): the simplification of an algebraic boundary, and its form with the singular
locus of $X$ contained in $W$.

\inputleanmodule{Hironaka.Ch4.MainTheoremIICorollary3}

\subsection*{The maximal-contact step of Lemma (CD.4) (p.~318)}

The paper's citation of Lemma~22 of Ch.~III for ``$z_{t+1}$ is a non-unit'' presupposes its
conclusion; these nodes prove the claim along the proof of Lemma~22.

\inputleanmodule{Hironaka.Ch4.MaximalContact}

\subsection*{Normal crossings on $\tilde X$ (p.~321)}

The paper's ``we see'' for the normal crossings of $\hat E(s)$ with $\tilde X$ and of the $\tilde E_j$ on
$\tilde X$: the local algebra, in general.

\inputleanmodule{Hironaka.Compat.NormalCrossingsPullback}

\subsection*{The exceptional divisors and $\tilde X$ (pp.~308, 321)}

\inputleanmodule{Hironaka.Compat.ExceptionalTransversal}

\subsection*{Repair of the second error: splitting and emptying the $\tilde E_j$ (pp.~308, 310, 321--323)}

The auxiliary datum $\mathfrak R^{N-1}_{\mathrm{II}}$ of \S\S3--4 need not be a resolution datum (the $\tilde E_j$ can be
reducible, and $\tilde E_i=\tilde X$ when $F=E_i$; see the paper's error on the subschemes $\tilde E_j$, pp.~308, 315, 321). The nodes below split each
$\tilde E_j$ into its connected components and replace each $\tilde E_j=\tilde X$ by the empty subscheme; Theorem
$\mathrm{II}_2^{N-1}$ is applied to the split datum.

\inputleanmodule{Hironaka.Compat.ResolvableData}

\inputleanmodule{Hironaka.Compat.DatumRefinement}

\subsection*{Presentations of a datum (gap filled)}

\inputleanmodule{Hironaka.Compat.DatumPresentationBase}

\subsection*{The completion preserves orders (Lemma (CD.5), p.~320)}

\inputleanmodule{Hironaka.Compat.CompletionOrder}

\inputleanmodule{Hironaka.Compat.CompletionStalk}
